\documentclass[12pt,a4paper]{report}
\usepackage[top=2.5cm, bottom=2.5cm, left=2.5cm, right=2.5cm]{geometry}
\usepackage[utf8]{inputenc}
\usepackage[T1]{fontenc}
\usepackage[francais,english]{babel}
\usepackage[french]{varioref}
\usepackage{fancyhdr}
\IfFileExists{lmodern.sty}{\usepackage{lmodern}}{}

\pagestyle{fancy}
\fancyhead[L]{}\fancyhead[C]{}\fancyhead[R]{}
\fancyfoot[L]{}
\fancyfoot[C]{\small Adresse du site: Route de Soliman, Technopole de Borj Cédria, BP 123, Hammam Chatt 1164, Ben Arous, Tunisie - Tél: (216) 79 32 67 63/90/91/92 - Fax: (216) 79 32 67 64 - Site web: www.istic.rnu.tn}
\fancyfoot[R]{}
\renewcommand{\headrulewidth}{0pt}
\renewcommand{\footrulewidth}{0.1pt}

\begin{document}
\newcommand{\HRule}{\rule{\linewidth}{0.5mm}}

\textbf{Abstract} \\
\HRule \\
This project addresses the challenge of knowledge management at Jade Advisory by
designing and implementing Ask-Jade, an intelligent question-answering system based
on a Retrieval-Augmented Generation (RAG) architecture. The system transforms the
organization's dispersed document library into a centralized and searchable knowledge
base, enabling consultants to retrieve relevant project information efficiently. The
ingestion pipeline relies on Docling, BGE-M3, and pgvector, orchestrated through
LlamaIndex. A document management interface allows administrators to ingest, version,
and monitor documents from a connected SharePoint instance. User and model management
interfaces enable account administration and language model configuration without
redeployment. Users can also select their preferred model directly from the chat
interface. Evaluation results confirm the approach's relevance, achieving a
hallucination rate below the target threshold alongside high faithfulness and
helpfulness scores. The system supports English, French, and Arabic, with full
answer traceability through source citations.

\smallskip
\textbf{Keywords:} Retrieval-Augmented Generation (RAG), LlamaIndex, BGE-M3,
pgvector, Docling, knowledge management, large language model, vector search,
multilingual processing.

\vskip2cm

\textbf{Résumé} \\
\HRule \\
Ce projet répond au défi de la gestion des connaissances au sein de Jade Advisory en
concevant et en implémentant Ask-Jade, un système de question-réponse basé sur une
architecture RAG. Le système transforme la bibliothèque de documents dispersés en
une base de connaissances centralisée et interrogeable. Le pipeline d'ingestion repose
sur Docling, BGE-M3 et pgvector, orchestrés via LlamaIndex. Une interface de gestion
documentaire permet d'ingérer, versionner et surveiller les documents depuis SharePoint.
La gestion des utilisateurs et des modèles permet d'administrer les comptes et de
configurer les modèles de langage sans redéploiement. Les utilisateurs peuvent
également choisir leur modèle depuis l'interface de chat. Les résultats confirment
la pertinence de l'approche, avec un taux d'hallucination inférieur au seuil cible
et des scores élevés en fidélité et en utilité. Le système prend en charge l'anglais,
le français et l'arabe, avec une traçabilité complète via des citations sources.

\smallskip
\textbf{Mots clés :} Génération Augmentée par Récupération (RAG), LlamaIndex,
BGE-M3, pgvector, Docling, gestion des connaissances, modèle de langage,
recherche vectorielle, traitement multilingue.

\end{document}